\documentclass[preview]{standalone}

\usepackage{amsmath}
\usepackage{amssymb}
\usepackage{stellar}
\usepackage{definitions}
\usepackage{bettelini}

\begin{document}

\id{probability-countable-spaces}
\genpage

\section{Countable probability spaces}

\begin{snippetdefinition}{countable-probability-space-definition}{Countable probability space}
    Let \(\Omega\neq\emptyset\) be a finite or countably infinite \set.
    A \emph{countable probability space} on \(\Omega\) is a
    \probabilityspace
    \[
        (\Omega,\powerset(\Omega),\mathbb P)
    \]
    where \(\mathbb P\colon\powerset(\Omega)\fromto[0,1]\) satisfies:
    \begin{enumerate}
        \item \(\mathbb P(\Omega)=1\);
        \item for every pairwise disjoint family
            \(\{A_n\}_{n\in\naturalnumbers^\exceptzero}
            \subseteq\powerset(\Omega)\),
            \[
                \mathbb P\left(\bigcup_{n=1}^\infty A_n\right)
                =
                \sum_{n=1}^\infty \mathbb P(A_n).
            \]
    \end{enumerate}
    The second condition is called \(\sigma\)-additivity.
\end{snippetdefinition}

\begin{snippet}{countable-additivity-motivation}
    Countable additivity is the property that makes infinite probabilistic
    models compatible with limits of \event[events]. For example, in an infinite
    sequence of independent trials with success probability \(p>0\), let
    \(D_n\) be the \event of obtaining at least one success in the first
    \(n\) trials. Then \(D_n\subseteq D_{n+1}\) and the finite computation
    gives
    \[
        \mathbb P(D_n)=1-(1-p)^n.
    \]
    The expected conclusion for
    \(D=\bigcup_{n=1}^\infty D_n\) is
    \[
        \mathbb P(D)=\lim_{n\to\infty}\mathbb P(D_n)=1.
    \]
    This passage from an increasing sequence of \event[events] to its limit is a
    consequence of \(\sigma\)-additivity.
\end{snippet}

\section{Event limits}

\begin{snippetdefinition}{monotone-event-sequence-definition}{Monotone sequences of events}
    Let \((\Omega,\mathcal F)\) be a \measurablespace, let
    \(\{A_n\}_{n\in\naturalnumbers^\exceptzero}\subseteq\mathcal F\), and
    let \(A\in\mathcal F\).
    \begin{enumerate}
        \item We write \(A_n\uparrow A\) if
            \(A_n\subseteq A_{n+1}\) for every
            \(n\in\naturalnumbers^\exceptzero\), and
            \[
                A=\bigcup_{n=1}^\infty A_n.
            \]
        \item We write \(A_n\downarrow A\) if
            \(A_n\supseteq A_{n+1}\) for every
            \(n\in\naturalnumbers^\exceptzero\), and
            \[
                A=\bigcap_{n=1}^\infty A_n.
            \]
    \end{enumerate}
\end{snippetdefinition}

\begin{snippettheorem}{probability-continuity-from-below}{Continuity from below}
    Let \((\Omega,\mathcal F,\mathbb P)\) be a \probabilityspace.
    If \(\{A_n\}_{n\in\naturalnumbers^\exceptzero}\subseteq\mathcal F\)
    and \(A_n\uparrow A\), then
    \[
        \mathbb P(A)=\lim_{n\to\infty}\mathbb P(A_n).
    \]
\end{snippettheorem}

\begin{snippetproof}{probability-continuity-from-below-proof}{probability-continuity-from-below}{Continuity from below}
    Define
    \[
        B_1=A_1,
        \qquad
        B_n=A_n\difference A_{n-1}\quad\text{for }n\geq2.
    \]
    The \event[events] \(B_n\) are pairwise disjoint and
    \[
        A=\bigcup_{n=1}^\infty B_n.
    \]
    By countable additivity,
    \[
        \mathbb P(A)
        =
        \sum_{n=1}^\infty \mathbb P(B_n)
        =
        \lim_{m\to\infty}\sum_{n=1}^m\mathbb P(B_n).
    \]
    Since \(\bigcup_{n=1}^m B_n=A_m\), finite additivity gives
    \[
        \sum_{n=1}^m\mathbb P(B_n)=\mathbb P(A_m).
    \]
    Hence \(\mathbb P(A)=\lim_{m\to\infty}\mathbb P(A_m)\).
\end{snippetproof}

\begin{snippettheorem}{probability-continuity-from-above}{Continuity from above}
    Let \((\Omega,\mathcal F,\mathbb P)\) be a \probabilityspace.
    If \(\{A_n\}_{n\in\naturalnumbers^\exceptzero}\subseteq\mathcal F\)
    and \(A_n\downarrow A\), then
    \[
        \mathbb P(A)=\lim_{n\to\infty}\mathbb P(A_n).
    \]
\end{snippettheorem}

\begin{snippetproof}{probability-continuity-from-above-proof}{probability-continuity-from-above}{Continuity from above}
    If \(A_n\downarrow A\), then \(\Omega\difference A_n\uparrow
    \Omega\difference A\). By \continuityfrombelow,
    \[
        \mathbb P(\Omega\difference A)
        =
        \lim_{n\to\infty}\mathbb P(\Omega\difference A_n).
    \]
    Using the complement rule,
    \[
        1-\mathbb P(A)
        =
        \lim_{n\to\infty}\left(1-\mathbb P(A_n)\right)
        =
        1-\lim_{n\to\infty}\mathbb P(A_n),
    \]
    and therefore \(\mathbb P(A)=\lim_{n\to\infty}\mathbb P(A_n)\).
\end{snippetproof}

\begin{snippettheorem}{countable-additivity-continuity-characterization}{Countable additivity and continuity at the empty event}
    Let \(\Omega\neq\emptyset\) be a \set, and let
    \[
        \mathbb P\colon\powerset(\Omega)\fromto[0,1]
    \]
    satisfy \(\mathbb P(\Omega)=1\) and finite additivity on disjoint
    \event[events]. Then the following conditions are equivalent:
    \begin{enumerate}
        \item \(\mathbb P\) is countably additive;
        \item for every decreasing sequence
            \(\{A_n\}_{n\in\naturalnumbers^\exceptzero}\subseteq
            \powerset(\Omega)\) such that \(A_n\downarrow\emptyset\),
            \[
                \lim_{n\to\infty}\mathbb P(A_n)=0.
            \]
    \end{enumerate}
\end{snippettheorem}

\begin{snippetproof}{countable-additivity-continuity-characterization-proof}{countable-additivity-continuity-characterization}{Countable additivity and continuity at the empty event}
    If \(\mathbb P\) is countably additive, the second condition follows from
    \continuityfromabove applied to \(A=\emptyset\).

    Conversely, assume the second condition. Let
    \(\{A_n\}_{n\in\naturalnumbers^\exceptzero}\subseteq\powerset(\Omega)\)
    be pairwise disjoint, and put
    \[
        A=\bigcup_{n=1}^\infty A_n,
        \qquad
        R_m=A\difference\bigcup_{n=1}^m A_n.
    \]
    Then \(R_m\downarrow\emptyset\). By finite additivity,
    \[
        \mathbb P(A)
        =
        \mathbb P\left(\bigcup_{n=1}^m A_n\right)+\mathbb P(R_m)
        =
        \sum_{n=1}^m\mathbb P(A_n)+\mathbb P(R_m).
    \]
    Taking the limit as \(m\to\infty\) and using
    \(\mathbb P(R_m)\to0\) gives
    \[
        \mathbb P(A)
        =
        \sum_{n=1}^\infty\mathbb P(A_n).
    \]
    Thus \(\mathbb P\) is countably additive.
\end{snippetproof}

\begin{snippetproposition}{sigma-subadditivity-probability}{Sigma-subadditivity}
    Let \((\Omega,\mathcal F,\mathbb P)\) be a \probabilityspace.
    For every family
    \(\{A_n\}_{n\in\naturalnumbers^\exceptzero}\subseteq\mathcal F\),
    \[
        \mathbb P\left(\bigcup_{n=1}^\infty A_n\right)
        \leq
        \sum_{n=1}^\infty\mathbb P(A_n).
    \]
\end{snippetproposition}

\begin{snippetproof}{sigma-subadditivity-probability-proof}{sigma-subadditivity-probability}{Sigma-subadditivity}
    Define
    \[
        B_1=A_1,
        \qquad
        B_n=A_n\difference\bigcup_{k=1}^{n-1}A_k
        \quad\text{for }n\geq2.
    \]
    The \event[events] \(B_n\) are pairwise disjoint,
    \(B_n\subseteq A_n\), and
    \[
        \bigcup_{n=1}^\infty B_n
        =
        \bigcup_{n=1}^\infty A_n.
    \]
    Hence, by countable additivity and monotonicity,
    \[
        \mathbb P\left(\bigcup_{n=1}^\infty A_n\right)
        =
        \sum_{n=1}^\infty\mathbb P(B_n)
        \leq
        \sum_{n=1}^\infty\mathbb P(A_n).
    \]
\end{snippetproof}

\section{Discrete densities}

\begin{snippetdefinition}{discrete-density-definition}{Discrete density}
    Let \(\Omega\neq\emptyset\) be a finite or countably infinite \set.
    A \emph{discrete density} on \(\Omega\) is a \function
    \[
        p\colon\Omega\fromto[0,1]
    \]
    such that
    \[
        \sum_{\omega\in\Omega}p(\omega)=1.
    \]
\end{snippetdefinition}

\begin{snippetproposition}{probability-density-associated-with-countable-space}{Density associated with a countable probability space}
    Let \((\Omega,\powerset(\Omega),\mathbb P)\) be a
    \countableprobabilityspace. The \function
    \[
        p\colon\Omega\fromto[0,1],
        \qquad
        p(\omega)=\mathbb P(\{\omega\})
    \]
    is a \discretedensity on \(\Omega\).
\end{snippetproposition}

\begin{snippetproof}{probability-density-associated-with-countable-space-proof}{probability-density-associated-with-countable-space}{Density associated with a countable probability space}
    Since \(\Omega\) is finite or countably infinite, it is the disjoint union
    of the singletons \(\{\omega\}\), with \(\omega\in\Omega\). Therefore
    \[
        \sum_{\omega\in\Omega}p(\omega)
        =
        \sum_{\omega\in\Omega}\mathbb P(\{\omega\})
        =
        \mathbb P(\Omega)
        =
        1.
    \]
\end{snippetproof}

\begin{snippettheorem}{probability-measure-from-discrete-density}{Probability measure induced by a discrete density}
    Let \(\Omega\neq\emptyset\) be a finite or countably infinite \set, and
    let \(p\colon\Omega\fromto[0,1]\) be a \discretedensity. Define
    \[
        \mathbb P\colon\powerset(\Omega)\fromto[0,1],
        \qquad
        \mathbb P(A)=\sum_{\omega\in A}p(\omega).
    \]
    Then \((\Omega,\powerset(\Omega),\mathbb P)\) is a
    \countableprobabilityspace.
    Conversely, if \((\Omega,\powerset(\Omega),\mathbb P)\) is a
    \countableprobabilityspace with associated \discretedensity \(p\), then
    \[
        \mathbb P(A)=\sum_{\omega\in A}p(\omega)
    \]
    for every \(A\subseteq\Omega\).
\end{snippettheorem}

\begin{snippetproof}{probability-measure-from-discrete-density-proof}{probability-measure-from-discrete-density}{Probability measure induced by a discrete density}
    First,
    \[
        \mathbb P(\Omega)=\sum_{\omega\in\Omega}p(\omega)=1.
    \]
    Let
    \(\{A_n\}_{n\in\naturalnumbers^\exceptzero}\subseteq\powerset(\Omega)\)
    be pairwise disjoint. Since all terms are nonnegative,
    \[
        \mathbb P\left(\bigcup_{n=1}^\infty A_n\right)
        =
        \sum_{\omega\in\bigcup_{n=1}^\infty A_n}p(\omega)
        =
        \sum_{n=1}^\infty\sum_{\omega\in A_n}p(\omega)
        =
        \sum_{n=1}^\infty\mathbb P(A_n).
    \]
    Hence \(\mathbb P\) is countably additive.

    Conversely, if \(\mathbb P\) is a \countableprobabilityspace, then every
    \(A\subseteq\Omega\) is the disjoint union of the singletons
    \(\{\omega\}\) with \(\omega\in A\). Therefore
    \[
        \mathbb P(A)
        =
        \sum_{\omega\in A}\mathbb P(\{\omega\})
        =
        \sum_{\omega\in A}p(\omega).
    \]
\end{snippetproof}

\begin{snippetcorollary}{discrete-density-determines-probability-measure}{A discrete density determines its probability measure}
    Let \(\Omega\neq\emptyset\) be a finite or countably infinite \set.
    Let \(\mathbb P_1\) and \(\mathbb P_2\) be
    \probabilitymeasure[probability measures] on
    \((\Omega,\powerset(\Omega))\), with associated \discretedensity[discrete
    densities] \(p_1\) and \(p_2\). Then
    \[
        \mathbb P_1=\mathbb P_2
        \quad\ifandonlyif\quad
        p_1=p_2.
    \]
\end{snippetcorollary}

\begin{snippetproof}{discrete-density-determines-probability-measure-proof}{discrete-density-determines-probability-measure}{A discrete density determines its probability measure}
    If \(\mathbb P_1=\mathbb P_2\), then for every \(\omega\in\Omega\),
    \[
        p_1(\omega)=\mathbb P_1(\{\omega\})
        =
        \mathbb P_2(\{\omega\})=p_2(\omega).
    \]
    Conversely, if \(p_1=p_2\), then for every \(A\subseteq\Omega\),
    \[
        \mathbb P_1(A)
        =
        \sum_{\omega\in A}p_1(\omega)
        =
        \sum_{\omega\in A}p_2(\omega)
        =
        \mathbb P_2(A).
    \]
\end{snippetproof}

\begin{snippetcorollary}{discrete-density-support-full-measure}{Support of a discrete density}
    Let \(\Omega\neq\emptyset\) be a finite or countably infinite \set, let
    \(p\) be a \discretedensity on \(\Omega\), and let \(\mathbb P\) be the
    induced \probabilitymeasure. Define
    \[
        S=\{\omega\in\Omega\suchthat p(\omega)>0\}.
    \]
    Then \(\mathbb P(S)=1\). Moreover, \(\mathbb P\) gives probability
    \(1\) to some finite \set if and only if \(S\) is finite.
\end{snippetcorollary}

\begin{snippetproof}{discrete-density-support-full-measure-proof}{discrete-density-support-full-measure}{Support of a discrete density}
    Since \(p(\omega)=0\) for every \(\omega\in\Omega\difference S\),
    \[
        \mathbb P(\Omega\difference S)
        =
        \sum_{\omega\in\Omega\difference S}p(\omega)
        =
        0.
    \]
    Hence \(\mathbb P(S)=1\).

    If \(S\) is finite, then \(S\) itself is a finite \set with
    \(\mathbb P(S)=1\). Conversely, suppose that \(F\subseteq\Omega\) is
    finite and \(\mathbb P(F)=1\). Then
    \(\mathbb P(\Omega\difference F)=0\), so no
    \(\omega\in\Omega\difference F\) can satisfy \(p(\omega)>0\). Therefore
    \(S\subseteq F\), and \(S\) is finite.
\end{snippetproof}

\section{Countable total probability}

\begin{snippettheorem}{countable-total-probability-law}{Countable law of total probability}
    Let \((\Omega,\mathcal F,\mathbb P)\) be a \probabilityspace.
    Let
    \[
        \{B_n\}_{n\in\naturalnumbers^\exceptzero}\subseteq\mathcal F
    \]
    be a \samplespacepartition of \(\Omega\), and assume
    \(\mathbb P(B_n)>0\) for every \(n\in\naturalnumbers^\exceptzero\).
    Then, for every \(A\in\mathcal F\),
    \[
        \mathbb P(A)
        =
        \sum_{n=1}^\infty\mathbb P(A\intersection B_n)
        =
        \sum_{n=1}^\infty\mathbb P(A\mid B_n)\mathbb P(B_n).
    \]
\end{snippettheorem}

\begin{snippetproof}{countable-total-probability-law-proof}{countable-total-probability-law}{Countable law of total probability}
    Since \(\{B_n\}_{n\in\naturalnumbers^\exceptzero}\) is a
    \samplespacepartition of \(\Omega\), the \event[events]
    \(A\intersection B_n\) are pairwise disjoint and
    \[
        A=A\intersection\Omega
        =
        \bigcup_{n=1}^\infty(A\intersection B_n).
    \]
    Countable additivity gives
    \[
        \mathbb P(A)
        =
        \sum_{n=1}^\infty\mathbb P(A\intersection B_n).
    \]
    Since \(\mathbb P(B_n)>0\), the definition of \conditionalprobability gives
    \[
        \mathbb P(A\intersection B_n)
        =
        \mathbb P(A\mid B_n)\mathbb P(B_n),
    \]
    and the formula follows.
\end{snippetproof}

\end{document}
